%%
%% spring-lecture04.tex
%% 
%% Made by Alex Nelson <pqnelson@gmail.com>
%% Login   <alex@lisp>
%% 
%% Started on  2026-04-09T10:40:40-0700
%% Last update 2026-04-09T10:40:40-0700
%% 

\lecture[Riemannian manifolds]

% START bit from previous lecture which fits more logically here
\begin{definition}[Riemannian metric]
Let $M$ be a smooth $n$-manifold. A \define{Riemannian Metric} on $M$
is a $(0,2)$-tensor $g$ on $M$ which is:
\begin{enumerate}
\item \textbf{Symmetric}: for each $x\in M$, for each $v,w\in T_{x}M$, $g(x)[v,w]=g(x)[w,v]$
\item \textbf{Nondegenerate}: for each $x\in M$, for each $v\in T_{x}M$,
  if we have for every $w\in T_{x}M$ that $g(x)[v,w]=0$, then $v=0$
\item \textbf{Positive}: for each $x\in M$ and for each $v\in T_{x}M$,
  we have $g(x)[v,v]\geq0$ and (by nondegeneracy) $g(x)[v,v]=0$ if only
  if $v=w$.
\end{enumerate}
If we relax the ``positive'' condition, then we get a semi-Riemannian
metric.
\end{definition}
% END bit from previous lecture which fits more logically here

\begin{definition}
A \define{Riemannian Manifold} $(M,g)$ consists of a smooth manifold
$M$ and a Riemannian metric $g$ on $M$.
\end{definition}

\begin{definition}[Local coefficients]
Let $(M,g)$ be a Riemannian manifold. Let
$(U,\varphi=(\xi^{1},\dots,\xi^{n}))\in\smoothstr{M}$ be a local
chart. Then the \define{Local Coefficients} of $g$ on $(U,\varphi)$ is
given by
\begin{equation}
g_{ij}(x):=g(x)\left[\frac{\partial}{\partial\xi^{i}}(x),\frac{\partial}{\partial\xi^{j}}(x)\right]
\end{equation}
for any $x\in U$, and for any $i,j=1,\dots,n=\dim(M)$. Observe
$g_{ij}\colon U\to\RR$ are all smooth functions.
\end{definition}

\begin{remark}
The Euclidean metric on $\RR^{n}$ is given by
\begin{equation}
g_{0}:=\sum_{i}(\D x_{i})^{2}.
\end{equation}
We can generalize this! Let $(M,g)$ be a Riemannian manifold, let
$(U,\varphi=(\xi^{1},\dots,\xi^{n}))\in\smoothstr{M}$ be a local
chart. Then for every $i=1,\dots,n$ the differential 1-form
$\D\xi^{i}\in\Omega^{1}(U)$ is given by
\begin{equation}
\D\xi^{i}(x)[v]=v_{i}
\end{equation}
where
\begin{equation}
v=\sum^{n}_{i=1}v_{i}\frac{\partial}{\partial\xi^{i}}(x).
\end{equation}
This is just the dual basis for the cotangent space. With this
notation, we have
\begin{equation}
g=\sum^{n}_{i,j=1}g_{ij}\,\D\xi^{i}\otimes\D\xi^{j}
\end{equation}
on $U$. This means
\begin{equation}
g(x)[v,w]=\sum^{n}_{i,j=1}g_{ij}(x)v_{i}w_{j}=\sum g(x)[\frac{\partial}{\partial\xi^{i}}(x),\frac{\partial}{\partial\xi^{j}}(x)]v_{i}w_{j}.
\end{equation}
This is completely formal, we will forget about it in practice.
\end{remark}

\begin{remark}
We should remember that second countable plus locally Euclidean
implies paracompactness.
\end{remark}

\begin{example}
$(\RR^{n},g_{0})$ is a Riemannian manifold with
\begin{equation}
g_{0}(x)[v,w]=\sum^{n}_{i}v_{i}w_{i}
\end{equation}
for any $x\in\RR^{n}$ and $v,w\in T_{x}\RR^{n}$. This $g_{0}$ is
called the \define{Standard Riemannian Metric} on $\RR^{n}$.
\end{example}

\begin{example}[Conformally equivalent metrics]
Given a Riemannian manifold $(M,g)$ and a smooth $f\colon M\to(0,+\infty)$,
then the $(0,2)$-tensor $fg$ on $M$ given by
\begin{equation}
(fg)(x)[v,w]=f(x)g(x)[v,w]
\end{equation}
is also a Riemannian metric on $M$. If we write $h=fg$ for this
metric, then $h$ and $g$ are said to be \define{Conformally Equivalent}.
If $f$ vanishes at finitely many points, then $h$ and $g$ are
\define{Weakly Conformally Equivalent}.
\end{example}

\begin{example}[Pullback metric]
Let $(N,h)$ be a Riemannian manifold. Let $M$ be a smooth manifold.
Let $f\colon M\to N$ be a smooth immersion. Then the $(0,2)$-tensor
$f^{*}h$ given by
\begin{equation}
(f^{*}h)(x)[v,w]:=h(f(x))[\D f(x)[v],\D f(x)[w]],
\end{equation}
for any $x\in M$ and $v,w\in T_{x}M$, is a Riemannian metric on
$M$ called the \define{Pullback Metric} of $h$ along $f$.
If $f$ is not an immersion, then you can get degeneracies. This
is a great way to construct metrics.
\end{example}

\begin{example}
All submanifolds $M$ of $\RR^{n}$ ar eembeddings, and this is one way
to get a matric by using the inclusion $i\colon M\into\RR^{n}$, $i(x)=x$.
Then inclusion $i$ is an embedding (since $M$ is a smooth submanifold
of $\RR^{n}$) so $g=i^{*}g_{0}$ is a Riemannian metric on $M$, meaning
\begin{equation}
g(x)[v,w]:=g_{0}(x)[v,w]
\end{equation}
for any $x\in M$ and $v,w\in T_{x}M\subset T_{x}\RR^{n}\iso\RR^{n}$.
\end{example}

\begin{question}
If we have a smooth submanifold $M$ of $\RR^{n}$, it is naturally a
Riemannian manifold with the pullback metric. If we have a Riemannian
manifold $(M,g)$, is it a submanifold of some $\RR^{n}$?
\end{question}

\begin{theorem}[Nash's embedding]
Let $(M,g)$ be a Riemannian manifold. Then there exists an embedding
$f\colon M\into\RR^{k}$ (for sufficiently large $k$) such that $f^{*}g_{0}=g$.
\end{theorem}

The currently known upper bounds on $k$ is: if $M$ is compact, then
$k\leq n^{2}$. If $M$ is noncompact, then $k\leq n^{3}$.

(The proof is long and complicated. It would take a week to discuss.)

\begin{question}
Is every smooth manifold a Riemannian manifold?
\end{question}

\begin{theorem}
Every smooth $n$-manifold admits at least one Riemannian metric.
\end{theorem}

\begin{proof}
Since $M$ is second-countable and locally Euclidean, we see $M$ is
paracompact. Since $M$ is paracompact, Hausdorff, and
second-countable, there exists a countable atlas
$\{(U_{k},\varphi_{k})\}_{k\in\NN}$ and a locally finite partition of
unity subordinated to the open cover $\{U_{k}\}_{k\in\NN}$ (i.e., a
family of smooth maps $\rho_{k}\colon U_{k}\to[0,1]$ where $k\in\NN$
such that
\begin{enumerate}
\item $\rho_{k}$ is compactly supported in $U_{k}$, i.e.,
  $\supp(\rho_{k})\subset U_{k}$ is compact;
\item for all $x\in M$, $\rho_{k}(x)\neq0$ for only finitely many $k\in\NN$
\item for all $x\in M$, $\sum_{k\in\NN}\rho_{k}(x)=1$.
\end{enumerate}
This is the notion of a partition of unity.) Then we define the
$(0,2)$-tensor $g$ on $M$ by
\begin{equation}
g(x)=\sum_{k\in\NN}\rho_{k}(x)(\varphi^{*}_{k}g_{0})(x).
\end{equation}
It's not hard to prove $g$ is a Riemannian metric, which proves the result.
\end{proof}

\begin{xca}
Work out the details of the proof.
\end{xca}

%% START stuff from next lecture placed here for logical coherence
\begin{definition}[Isometry]
Let $(M,g)$ and $(N,h)$ be Riemannian manifolds. A diffeomorphism
$f\colon M\to N$ is called an \define{Isometry} if $f^{*}h=g$. In
other words, for all $x\in M$ and for all $v,w\in T_{x}M$ we have
\begin{equation}
h(f(x))[\D f(x)[v],\D f(x)[w]]=g(x)[v,w].
\end{equation}
\end{definition}

\begin{node}
Just to review, we have a little table of isomorphisms
\begin{center}
\begin{tabular}{c|c}
Category & Isomorphism \\\hline
Topological Manifolds & Homeomorphisms\\
Smooth Manifolds & Diffeomorphisms\\
Riemannian Manifolds & Isometries
\end{tabular}
\end{center}
\end{node}
%% END stuff from next lecture placed here for logical coherence